\newpage
\pagestyle{fancy}

\section{Premessa}

Il presente documento costituisce la relazione tecnica e specialistica del progetto focalizzato sugli interventi necessari per adeguare gli impianti elettrici esistenti in modo da garantire l’installazione di una stazione di ricarica per i veicoli elettrici nel box condominiale di proprietà di \{nomeCommittente\} \{cognomeCommittente\} presso \{indirizzoCommittente\}, in conformità alle normative vigenti. \\
Il presente documento è redatto in conformità al Decreto Ministeriale 22/01/2008 N°37, il quale stabilisce che la progettazione degli impianti elettrici deve essere obbligatoriamente affidata a un professionista iscritto all'Albo professionale e competente nel settore. \\
La relazione illustra puntualmente e nel dettaglio tutti gli aspetti esaminati e le verifiche analitiche effettuate in sede di progettazione, fornendo un quadro completo delle opere da eseguire, corredato di documenti annessi e richiamando le normative tecniche, i decreti e le circolari che dovranno essere rigorosamente seguiti come parte integrante del progetto. \\
I materiali e i dispositivi utilizzati che comporranno l’impianto oggetto del seguente intervento devono rispettare, nei contesti pertinenti, le prescrizioni contenute nelle norme di cui al capitolo 2, compresi eventuali varianti, aggiornamenti ed estensioni emanate posteriormente dargli organismi di normazione citati di seguito. In particolare, \textbf{i materiali e gli apparecchi utilizzati devono essere conformi alle normative CEI e CEI-UNEL}, in modo tale da essere adatti all’ambiente in cui sono installati con resistenza alle sollecitazioni meccaniche, corrosione, variazioni termiche, umidità e acque meteoriche alle quali possono essere esposti durante il funzionamento. \\
Quanto precedentemente citato comprende la fornitura e installazione di tutti i componenti, anche quelli che non sono espressamente citati ma necessari per un corretto funzionamento dell’impianto oggetto del seguente intervento, conforme, quindi, agli standard di settore, obbligatori e volontari. \\
Tutti i materiali e i dispositivi, inoltre, per i quali è prevista la concessione del marchio di qualità sono muniti del contrassegno IMQ o CE. \\
\textbf{Al momento della consegna, l’impianto dovrà rispettare tutte le normative e leggi applicabili.}


\subsection{Tecnico Progettista}
\begin{center}
\begin{tabular}{|p{6.5cm}|p{7cm}|}
\hline
\textbf{Nome e Cognome} & Giuseppe CARDELLO \\
\hline
\textbf{Qualifica} & Ingegnere \\
\hline
\textbf{Codice Fiscale} & CRDGPP96T19C349U \\
\hline
\textbf{P.IVA} & 13478910964 \\
\hline
\textbf{Data di nascita} & 19/12/1996 \\
\hline
\textbf{Indirizzo di residenza} & Via Walter Tobagi, 7 \\
\hline
\textbf{CAP -- Comune} & 20143 - Milano (MI) \\
\hline
\textbf{Albo} & Ingegneri \\
\hline
\textbf{Provincia di iscrizione} & Milano \\
\hline
\textbf{Numero di iscrizione} & A 34705 \\
\hline
\textbf{Telefono} & +39 3407580346 \\
\hline
\textbf{E-mail} & GCardello96@gmail.com \\
\hline
\end{tabular}
\end{center}